\section{Introduction}\label{sec:introduction}

A tangent flow describes the geometry of a mean curvature flow at a
singularity after magnifying space and time. Compactness produces such
limits along subsequences; uniqueness asks whether every magnification
sequence gives the same limiting motion, including its placement in the
fixed ambient frame. We prove this uniqueness at the first singular time of
every smooth closed embedded surface in three-dimensional Euclidean space.

Let $M_0\subset\R^3$ be a smooth compact connected embedded surface
without boundary, and let $(M_t)_{0\le t<T}$ be its maximal smooth mean
curvature flow. The normal velocity is the trace mean curvature vector.
For a singular point $(x_0,T)$ and a scale $\lambda>0$, define
\begin{equation}\label{intro:rescalings}
 M_s^\lambda=\lambda\bigl(M_{T+\lambda^{-2}s}-x_0\bigr),\qquad s<0,
\end{equation}
whenever the original time belongs to $[0,T)$.  Write
$\mu_s^\lambda=\mathcal H^2\!\llcorner M_s^\lambda$ for its
multiplicity-one area measure.  The center $x_0$ and the ambient
coordinates stay fixed throughout.

We use the following precise convention for the weak-flow consequence.
A \emph{backward integral Brakke tangent} is a weakly measurable Radon
family $(\mu_s)_{s<0}$ with the following properties.  For almost every
$s$, $\mu_s$ is the weight of an integral $2$-varifold $V_s$ with
measurable generalized mean curvature $\mathbf H_\mu$ and first variation
$\delta V_s=-\mathbf H_\mu\mu_s$ on compactly supported vector fields.
The family satisfies the local integrability and all-endpoint conditions
\eqref{geo:representative-integrability}--\eqref{geo:all-endpoint-brakke},
and some $\lambda_i\to\infty$ gives local spacetime weight convergence
\[
 \mu_s^{\lambda_i}\,ds\ \rightharpoonup\ \mu_s\,ds
 \quad\hbox{on }\mathbb R^3\times(-\infty,0).
\]
The all-endpoint condition allows compact tests nonzero at either
endpoint and requires the inequality for every time pair.  This is the
endpoint/test clause used in \cite[Section~2.1]{BK2024}; the weak-flow
framework is due to Brakke \cite{Brakke1978}.  An extraction with locally
weak slice-weight convergence for almost every time and uniform local
mass bounds on compact time intervals implies the displayed premise.
The area-ratio bounds here supply that domination.  In particular, an
extraction that also has local varifold convergence for almost every
time, or has the stronger slice-weight convergence at every time, is
included when its representative satisfies the all-endpoint condition.

\begin{theorem}\label{thm:main}
Let $M_0\subset\R^3$ be a smooth compact connected embedded surface
without boundary, and let $(M_t)_{0\le t<T}$ be its maximal smooth
mean curvature flow.  At every singular point $(x_0,T)$ at its first
singular time there is a smooth properly embedded self-shrinker $S$
such that, as $\lambda\to\infty$, the full family
\eqref{intro:rescalings} converges locally smoothly in space and with
multiplicity one to $\sqrt{-s}\,S$, uniformly for $s$ in compact
subintervals of $(-\infty,0)$.  The convergence is in the original
ambient coordinates, including the axes and position of the limit.
Every backward integral Brakke tangent in the stated convention
therefore equals the area-measure flow of $\sqrt{-s}\,S$ at every
negative time.
\end{theorem}

There is no mean-convexity assumption and no prescribed tangent model.
The entropy and genus bounds used in the proof are automatic for the
compact initial surface. The theorem concerns a fixed center and the
smooth evolution before the first singular time; moving-center
rescalings and continuations through a singularity are different questions.

\subsection{Geometric background and preceding uniqueness results}

Huisken's monotonicity formula identifies self-shrinking solutions as
the natural models for parabolic blowups \cite{Huisken1990}. In the
normalization used here, a self-shrinker $S$ satisfies
$\mathbf H+X^\perp/2=0$ and generates the motion $\sqrt{-s}\,S$.
It is a critical point of the Gaussian area
\[
 F(S)=\int_S e^{-|X|^2/4}\,d\mu.
\]
Ilmanen's work establishes smoothness of tangent supports at the first
singular time of closed embedded surface flows in $\R^3$
\cite[Theorems~1 and~2]{Ilmanen1995}. Smooth support alone
does not settle multiplicity or imply smooth convergence of the
approximating flows. Lu Wang proved that each end of a noncompact
properly embedded self-shrinker in $\R^3$ of finite topology approaches
a regular cone under dilations, smoothly on compact sets away from the
vertex, or a round cylinder under translations, smoothly on compact
sets \cite[Theorem~1.1]{Wang2016}. The
multiplicity-one theory of Bamler--Kleiner
\cite{BK2024}, together with their tangent structure and end theorem,
directly gives the geometric input used here: a smooth properly embedded tangent
section with finitely many separated conical and round cylindrical
ends. \Cref{geo:tangent} states that input and checks its
vanishing-density-drop hypothesis. The passage to smooth local
convergence uses the local Brakke-regularity consequence recorded by
Schulze \cite[Lemma~2.1]{Schulze2014}; see also White \cite{White2005}.

Several major classes of tangent models already have uniqueness
theorems. Schulze treats a smooth compact embedded multiplicity-one
shrinker, in arbitrary dimension and codimension
\cite[Corollary~1.2]{Schulze2014}. Colding--Minicozzi prove uniqueness
for multiplicity-one round cylindrical hypersurface tangents, including
their axes \cite[Theorem~0.2]{CM2015}. Chodosh--Schulze establish
uniqueness for multiplicity-one asymptotically conical hypersurface
tangents \cite[Theorem~1.1]{CS2021}. These conclusions determine the ambient
placement of the limiting model, including its axis in the cylindrical case.
In the same first-time surface setting, Bernstein--Wang also
obtained a uniqueness statement modulo rotations for planar,
spherical, and cylindrical models, allowing higher multiplicities
\cite[Theorem~1.1]{BW2015}.  Recent alternative approaches in the
cylindrical class include Ghosh's comparison-flow method
\cite[Theorem~1.1 and Sections~4 and~6]{Ghosh2026} and the quantitative
PDE--ODI method of Bamler--Lai
\cite[Corollary~5.97]{BamlerLai2026}.

A shrinker with cylindrical ends need not have been identified as an
exact round cylinder. The finite-end structure theorem therefore
leaves a stationary analysis involving all its conical and cylindrical
ends at once. This distinction is also visible in Haslhofer's
formulation of the cylinder-rigidity and tangent-uniqueness problems
\cite[Section~5.1, Problem~5.1 and Conjecture~5.2]{Haslhofer2025}.
\Cref{thm:main} establishes the negative-time part of that fixed-center
uniqueness assertion at the first singular time, without using cylinder rigidity.

\subsection{The stationary obstruction}

Gradient inequalities connect small stationary defect to finite length
of the rescaled flow.  Simon developed this analytic convergence
method for nonlinear geometric equations \cite{Simon1983}; its
application to compact tangent flows is explicit in
\cite[Section~3]{Schulze2014}. The Gaussian compactness, finite-dimensional
kernel augmentation, and analytic reduction used in this approach have
close predecessors in Chodosh--Schulze
\cite[Sections~3.4 and~4]{CS2021}. For the present end geometry, the
obstacle is a Jacobi field that grows quadratically along a cylindrical
end. It changes the type of the end. Although Gaussian area is analytic
on suitable spaces of decaying graphs, that fact does not establish
analyticity in the opening directions. We instead retain finite Taylor jets.

The relevant scale is already visible in the circular comparison.
If $\beta<0$ is its quadratic opening coefficient, then
$R_c(z)^2=2(1+\beta z^2)$ reaches zero at
$z=(-\beta)^{-1/2}$.  The Gaussian factor there is
$\exp(-1/(4|\beta|))$.  Thus the geometry changes by order one at a
distance where the action records effects smaller than every power of
the opening.  The stationary argument compares those effects while
keeping only finite Taylor information at each step.

Finite-order obstruction estimates have a close precedent in Zhu's
study of exact product shrinkers with a closed simply nonintegrable
cross-section \cite[Section~1.3]{Zhu2024}. His finite-dimensional model
eliminates variables complementary to the kernel and derives a power
bound from a uniform leading obstruction at finite order. For the exact
products, the geometric argument establishes a second-order obstruction
in the Jacobi directions complementary to rotations. The argument below
obtains its required finite order by ruling out formal curves of critical
points with nonzero openings and applying real-algebraic certificates.

We choose compactly supported linear observations of a graph and solve
the stationary equation subject to their prescribed values. The
constraints introduce finitely many multipliers supported on the
observation core; outside that core the resulting surfaces are exact
shrinkers. Their parameters split as $(b,x)$. The analytic family
at $x=0$ preserves end types, while $x$ has one quadratic opening
coordinate for each cylindrical end. Successive linear solves define
a formal Gaussian action $\mathcal F(b,x)$ whose coefficients are
analytic in $b$. A formal critical arc is a tuple of real formal
power series $(b(s),x(s))$, with zero constant terms, satisfying
$\nabla\mathcal F(b(s),x(s))=0$ coefficient by coefficient.

The central stationary assertion is that every such arc has
$x(s)\equiv0$. Otherwise, two normalized stationary surfaces can be
constructed on complex contours passing on opposite sides of the zero
radius of a circular comparison profile. These are complex coefficients
on real spatial parameter domains; an arbitrary smooth graph is never
continued to a complex spatial argument. The two actions have a common
formal expansion with Gevrey bounds: factorial control of coefficients
and of the remainder at every order. A finite-dimensional implicit
replacement of the formal arc and a sector argument make the difference
of each action from the base critical value, as well as its gradient,
exponentially small at its own selected parameters.
Taylor transport compares both actions at common parameters with a
strict exponential margin.

The relation between Gevrey expansions and exponentially small
sectorial differences is classical; Malgrange--Ramis describe the
corresponding flatness and sectorial uniqueness principles
\cite[Sections~1.3--1.6]{MR1992}. Nonlinear Stokes theory for analytic
ordinary differential equations provides a related asymptotic framework
\cite{Costin1998}. Here the required estimates are proved for the
stationary surface equations, including their compact constraints,
prescribed contours, and endpoint errors.

The competing lower bound comes from the collapsing end. Height
coordinates reduce its exterior limit to the radial translator
equation. This is the same scalar equation whose real asymptotics
are analyzed by Clutterbuck--Schn\"urer--Schulze
\cite[Section~2]{CSS2007}. We prove that its complex lateral solutions
have a nonzero Stokes difference and compute the Gaussian action jump
by boundary momentum. A second-order mean calculation determines the
exponential phase to constant order. Distinct negative-power Laurent
principal parts separate at a generic nearby angle; ends with the same
entire principal part have positive limiting relative weights. Their
contributions cannot cancel. This contradicts the action upper bound,
including for even-order real arcs evaluated at complex parameters.

\subsection{Finite jets, localization, and convergence}

The formal real Nullstellensatz, in the arc formulation of
Aschenbrenner--Srhir \cite[Corollary~4.8]{AS2024}, turns the exclusion of opening arcs into
finite real-radical certificates for the gradient ideal. Artin's
approximation theorem \cite[Theorem~1.2]{Artin1968} converts the
required finite quotients into analytic certificates. A single exponent
then controls the openings by the gradient of every sufficiently
high fixed truncation of $\mathcal F$. We choose one such truncation
and apply the ordinary finite-dimensional analytic gradient inequality
\cite[Section~IV.9]{Lojasiewicz1993}; convergence of the full opening
series is unnecessary.

A shifted Gaussian inverse compares a finite stationary Taylor graph
with an actual flow slice. This comparison retains both the compact
multiplier and the spatial cutoff, including when there are no opening
variables. It yields a localized gradient inequality with the Gaussian
tail measured at the scale of the gradient norm. Elliptic improvement,
graphical pseudolocality \cite[Theorem~1.5]{INS2019}, and positive-time
graphical smoothing \cite{EckerHuisken1991} propagate complete graphs
with an increasing radius. Choosing each radius from all preceding
energy drops makes the localization errors summable. Finite rescaled
length then closes the compact-core bootstrap and gives convergence
at every intermediate rescaled time in the original ambient frame.

\subsection{Organization}

\Cref{sec:geometry} fixes the geometric inputs and normalization.
\Cref{sec:linear} constructs the normalized family, its formal action,
and the real finite-Taylor comparison from Gaussian and infinite-end
estimates.  \Cref{sec:arrays} introduces the lateral contours, proves
the finite-path gauge and inverse estimates, and constructs the action
arrays and their exponential upper bound.  \Cref{sec:jump,sec:phase} compute the nonzero jump,
determine its phase, and exclude a critical opening arc.
\Cref{sec:finite-jet} derives the finite-power inequality.
\Cref{sec:localized} proves the localized gradient inequality, and
\Cref{sec:flow} uses it to obtain finite rescaled length, full smooth
convergence, and equality of every negative-time tangent measure.
